\subsection{RQ5: When does sequential RL become necessary?}
\label{sec:rq5}

In the frozen environment a move carries traffic in the interval in which it is
made, so its first-interval effect is in the immediate reward. We remove that
property with the delayed-activation variant ($L=1$): a request is charged when
it is made and takes effect one interval later; everything else, including
the observation of pending requests, is unchanged
(Table~\ref{tab:delay}).

The ranking reverses. On the same roots, the bandit falls from
\DelayBanditLZero{} without delay to \DelayBandit{} with it, because its target
now contains the cost of a move but none of its benefit. PPO is essentially
unaffected (\DelayPPOLZero{} $\rightarrow$ \DelayPPO), so under delay PPO beats
the bandit by \DelayPPOVs{} [\DelayPPOVsLo, \DelayPPOVsHi] (better on
\DelayPPORootsBetter{} of \DelayPPORoots{} roots). The masked Q-learner with $\gamma=0.9$, whose target includes the discounted
value of the next state and hence the delayed effect, partly recovers:
\DelayQNine{} under delay against \DelayQNineLZero{} without, above the bandit
on \DelayQNineRootsBetter{} of \DelayQNineRoots{} roots (mean difference
\DelayQNineVs) but below PPO.\ifcsname DelayQFive\endcsname{} With $\gamma=0.5$
(an effective horizon of about two intervals) it reaches \DelayQFive{}
(\DelayQFiveLZero{} without delay; above the bandit on
\DelayQFiveRootsBetter{} of \DelayQFiveRoots{} roots).\fi

The non-learning controllers degrade too, because they react to a state that
the delayed move no longer matches (MILP-track \DelayRefmilptrackLZero{}
$\rightarrow$ \DelayRefmilptrack, greedy \DelayRefgreedyLZero{} $\rightarrow$
\DelayRefgreedy); MILP-track remains the strongest policy, which is a
reminder that ``sequential RL becomes necessary'' is a statement about
learners whose targets omit the delayed effect, not about optimization-based
control.

\begin{table}[t]
\centering\small
\caption{Delayed TE activation ($L=1$), test return. Learners: roots 42 and 314159
(``No delay'' on the same roots); non-learning references: same test episodes.
``vs.\ bandit'': paired difference under delay.}
\label{tab:delay}
\setlength{\tabcolsep}{3.5pt}
\begin{tabular}{@{}lrrrrr@{}}
\toprule
Policy & Roots & No delay & $L=1$ & vs.\ bandit & Rer./h \\
\midrule
\tabinput{tables/delay_results}
\bottomrule
\end{tabular}
\end{table}
